% sec-01-abstract.tex - Abstract and Keywords (draft scaffold)
\begin{abstract}
\noindent
Phase 1 oncology trials advance only when a sequence of large regulatory and
clinical documents is authored, reviewed, and approved. This paper shows how a
repository based large language model (LLM), driven by a single master prompt that
first writes and then executes a schedule of sub-prompts, can hasten the entire
Phase 1 process by generating every relevant large document through an auditable
mermaid to draft to full to final pipeline.

\draftinstr{Write a 200-260 word abstract covering: (1) the patient-centered
motivation (enrolled PDAC patients cannot wait for administrative slowdowns); (2)
the method (repository LLM, single prompt, sub-prompt schedule, real-time
auto-commit/auto-PR monitorability); (3) the six ACCELERATION document targets;
(4) the verification stance (mermaid grounding, name-matching, GitHub
monitorability); (5) the benefit-greater-than-risk conclusion and admin/prep time
compression. State paper v1.0 and repository v4.2.0. Source the motivation and
method from \texttt{inputs/llm-adoption/main.tex}, the document targets from
\texttt{research/document-types/*}, and the workflow from
\texttt{research/industry-workflow/*}.}
\end{abstract}

\draftinstr{Add the Keywords line (\textbackslash keywords): Pancreatic ductal
adenocarcinoma; Phase 1 clinical trial; large language model; document generation;
repository based LLM; sub-prompt schedule; mermaid to LaTeX; clinical trial
acceleration; daraxonrasib; verification before generation.}
\keywords{[draft placeholder; rendered fully in the full and final stages]}
